% Generated draft: requires mathematical and source review.
\documentclass[11pt]{article}
\usepackage[margin=25mm]{geometry}
\usepackage{fontspec}
\setmainfont{Latin Modern Roman}
\usepackage{amsmath,amssymb,mathtools,unicode-math}
\setmathfont{Latin Modern Math}
\usepackage{enumitem,needspace,xurl,hyperref}
\setlength{\parindent}{0pt}
\setlength{\parskip}{5pt}
\setlength{\emergencystretch}{3em}
\newcommand{\N}{\mathbb{N}}
\newcommand{\Z}{\mathbb{Z}}
\newcommand{\Q}{\mathbb{Q}}
\newcommand{\R}{\mathbb{R}}
\newcommand{\C}{\mathbb{C}}
\newcommand{\F}{\mathbb{F}}
\newcommand{\A}{\mathbb{A}}
\newcommand{\PP}{\mathbb{P}}
\newcommand{\E}{\mathbb{E}}
\newcommand{\eps}{\varepsilon}
\newcommand{\dd}{\,\mathrm{d}}
\begin{document}
\section*{UnsolvedMath ID: 30006572}
\textbf{Source problem code:} catalogue entry
\subsection{Definitions and mathematical statement}
For an integer \(d\geq 2\), set
\[
\Gamma_d=\operatorname{PSL}_{2d-1}(\Z)
=\operatorname{SL}_{2d-1}(\Z)/Z(\operatorname{SL}_{2d-1}(\Z)).
\]
Let \(\lambda_d\) be the left regular representation of \(\Gamma_d\) on \(\ell^2(\Gamma_d)\), defined by \((\lambda_d(g)\xi)(h)=\xi(g^{-1}h)\). Its group von Neumann algebra is
\[
M_d=L(\Gamma_d)=\{\lambda_d(g):g\in\Gamma_d\}'',
\]
with canonical tracial state \(\tau_d(x)=\langle x\delta_e,\delta_e\rangle\), where \(\delta_e\) is the unit vector at the identity. For \(1\leq p<\infty\), let \(L^p(M_d,\tau_d)\) be the completion of \(M_d\) for
\[
\|x\|_p=\bigl(\tau_d(|x|^p)\bigr)^{1/p},
\]
equipped with its canonical noncommutative \(L^p\) operator-space structure. If the exponent \(p=\infty\) is admitted, use \(L^\infty(M_d,\tau_d)=M_d\) with its usual operator-space structure.

An operator space \(X\) has the completely bounded approximation property if there are a constant \(C<\infty\) and a net of finite-rank linear maps \(T_\alpha:X\to X\) such that
\[
\sup_\alpha\|T_\alpha\|_{\mathrm{cb}}\leq C,
\qquad
\lim_\alpha\|T_\alpha x-x\|_X=0
\quad\text{for every }x\in X,
\]
where
\[
\|T\|_{\mathrm{cb}}
=\sup_{n\geq1}\|\operatorname{id}_{M_n}\otimes T:M_n(X)\to M_n(X)\|.
\]
In particular, when \(X=M_d\), this definition uses pointwise convergence in operator norm.

\textbf{Question, with unresolved exponent range.} For each integer \(d\geq2\), is there an admissible exponent \(p\neq2\) such that \(L^p(M_d,\tau_d)\) has the completely bounded approximation property? If \(\mathcal P\) denotes the intended exponent range, the natural per-dimension reading is
\[
\forall d\in\N,\ d\geq2,\quad
\exists p\in\mathcal P\setminus\{2\}:\quad
L^p(M_d,\tau_d)\text{ has the completely bounded approximation property}.
\]
The supplied record does not specify \(\mathcal P\), including whether either endpoint is allowed. The definitions above cover the conventional Banach range \(1\leq p\leq\infty\); they do not determine the intended range. On the displayed reading, \(p\), the approximation bound, and the approximating maps may depend on \(d\). A requirement of one exponent working for every \(d\) would be a different quantifier order and is not explicit in the source.

The original excerpt begins with a reference to an omitted hypothesis being ``essential.'' Its supplied assessment also refers to a known result ``away from 2'' without stating that result or its exponent bounds. Consequently, the surrounding hypothesis and the precise intended extension remain unresolved. The following Kakeya expression in the excerpt has an undefined function \(f_d\) and is not sufficiently specified to form an additional mathematical statement here.
\subsection{Short English statement}
For each integer \(d\geq2\), the question asks whether the noncommutative \(L^p\)-space of the group von Neumann algebra of \(\operatorname{PSL}_{2d-1}(\Z)\) admits uniformly completely bounded finite-rank approximations to the identity for some \(p\neq2\). The intended exponent range and omitted surrounding hypothesis require clarification. The supplied record labels the problem open; that status is not independently verified here.
\subsection{Sources}
\begin{itemize}
\item Supplied local registry record: \texttt{data/unsolvedmath/problems.json}, permanent UnsolvedMath ID \(30006572\), exact source problem code \texttt{catalogue entry}, title: \(L^p\) Approximation for Higher-Rank Group Factors. Collection: Oberwolfach Reports Open Problems. No individual proposer is identified; the supplied proposed year is 2026.
\item Supplied report citation: \emph{Noncommutative Harmonic Analysis and Quantum Information} (2026), Oberwolfach Reports, pp.\ 217--275, DOI \texttt{10.4171/owr/2026/4}. Source URL: \url{https://doi.org/10.4171/owr/2026/4}. The record describes Oberwolfach Reports as documenting mathematical workshops and their open-problem discussions; no precise problem page or speaker attribution is supplied.
\item Supplied literature source: ``Oberwolfach Report 4/2026,'' \url{https://doi.org/10.4171/owr/2026/4}. Literature checked at: \texttt{2026-08-22}. This is supplied dated curation, not an independently verified assessment. It labels the question open and describes a known result only as being away from \(2\), without giving its precise statement or bounds.
\item The registry labels the statement \texttt{corrected\_verified} and says it was reconstructed and checked against the official Oberwolfach report. These are supplied provenance claims; the report has not been independently consulted for this draft. The supplied background requests independent verification of time-sensitive status and cited claims.
\end{itemize}
\end{document}
